% =============================================================================
% Conclusion
% =============================================================================
\section{Conclusion}\label{sec:conclusion}

We have studied the practical identifiability of viscoelastic shell parameters
from resonant frequencies in the spirit of Kac's spectral question.  For the
oblate Ritz flexural model used here, the inverse problem based on the
parameter triple $(a,c,E)$ transitions smoothly from poorly conditioned at the
near-spherical limit to well conditioned at the canonical oblate operating
point.  In the context of the Browntone programme, that makes Paper~8 less a
side study than a foundation stone: it addresses the inverse question on which
the diagnostic ambitions of Papers~1--7 depend.

The main findings are:
\begin{enumerate}
  \item within the Ritz model, the condition number drops from a near-spherical
  floor value $\kappa_\mathrm{floor}\approx269$ (five modes) to
  $\kappa = 69.4$ (canonical oblate), purely through geometric eccentricity;
  \item the equivalent-sphere model is algebraically rank deficient
  ($\kappa = 1.37 \times 10^{10}$), illustrating how model-induced degeneracy
  dwarfs the geometric effect;
  \item the power law $\kappa\sim C\,\varepsilon^{-\alpha}$ has theoretical
  exponent $\alpha=2$ and empirical estimates $\hat\alpha\in[2.16,\,2.66]$;
  \item round-trip inversion recovers $(a,c,E)$ to machine precision
  ($<10^{-12}$), confirming internal consistency of the forward model and
  solver;
  \item at $1\%$ frequency noise (five modes), the Cram\'er--Rao bounds are
  approximately $\pm 22.5\%$ for $a$, $\pm 6.0\%$ for $c$, and $\pm 30.8\%$
  for $E$, comparable to existing elastographic measurement precision;
  \item the condition number is stable across mode counts
  $\mathcal{N}=3$--$9$ ($\kappa\in[69,\,94]$), confirming that the
  identifiability improvement is robust and not an artefact of the Ritz
  truncation;
  \item three-mode inversion is vulnerable to local minima, whereas five-mode
  inversion is robust across the tested initial guesses;
  \item a ripe-watermelon case remains well conditioned, with
  $\kappa \approx 128$.
\end{enumerate}

These results suggest a practical answer to the question posed in the title.
Within the present model, one can hear the shape of an oblate organ in the
limited but useful sense of practical local identifiability.  One cannot do so
reliably at the spherical limit.  That distinction is precisely what determines
whether a measured resonance is merely interesting or diagnostically useful.

The programme-level implication is straightforward.  If flexural resonances can
be measured reproducibly, and if the relevant parameters are identifiable, then
the spectrum becomes a non-invasive elastography tool.  For abdominal
resonance~\cite{BrowntoneCompanion}, this means that a measured modal set could
in principle constrain effective cavity geometry and wall stiffness rather than
serving only as a forward prediction.  For the bladder
problem~\cite{BrowntoneBladder}, it means that fill-dependent resonance could,
at least in principle, support inference of changing wall properties.  For
watermelon ripeness~\cite{BrowntoneWatermelon}, it explains why a tap-tone can
carry recoverable material information rather than folklore alone.

The clearest practical recommendation is equally straightforward: measure at
least five low-order flexural resonances if the goal is parameter inversion
rather than mere modal description.  Fewer modes may suffice for local tests,
but they do not provide the same robustness.

The most interesting open mathematical target is the asymptotic law
\begin{equation}
  \kappa \sim C\,\varepsilon^{-\alpha}
  \qquad \text{as } \varepsilon \to 0^+,
\end{equation}
which is currently supported numerically (Table~\ref{tab:power_law}) and
formalised in Conjecture~\ref{conj:power_law}.  The theoretical prediction
$\alpha=2$ and the fitted prefactor $\hat C\approx24$--$26$ provide
quantitative design rules, but a rigorous proof would give the programme its
cleanest bridge between geometry and inverse conditioning: not merely that
oblateness helps, but exactly how fast identifiability is restored as one moves
away from the spherical limit.  That would be the crown jewel.  For the moment,
the jewel remains in the workshop.

Future work should therefore proceed in two directions.  First, the
Ritz approximation should be replaced by a full finite-element oblate shell
model to eliminate truncation error.  Second, the inversion should be tested
experimentally on phantoms and biological specimens.  Together, those steps
would move the Browntone programme from ``the resonances can be predicted'' to
``the material parameters can be recovered'', which is the more consequential
claim.
